\documentclass[11pt,a4paper]{article}
\usepackage[margin=2.2cm]{geometry}
\usepackage{lmodern}
\usepackage[T1]{fontenc}
\usepackage{xcolor,booktabs,tabularx,enumitem,titlesec}
\usepackage[most]{tcolorbox}
\usepackage[colorlinks=true,urlcolor=blue!60!black,linkcolor=black]{hyperref}
\usepackage{fontawesome5}
\usepackage{xurl}

\definecolor{accent}{HTML}{2A78D6}
\definecolor{warn}{HTML}{D1495B}
\definecolor{ok}{HTML}{2E8B57}
\setlength{\parindent}{0pt}
\setlength{\parskip}{4pt}
\setlist{itemsep=2pt,topsep=3pt}
\titleformat{\section}{\Large\bfseries\color{accent}}{}{0pt}{}
\titleformat{\subsection}{\large\bfseries}{}{0pt}{}
\titlespacing*{\section}{0pt}{14pt}{6pt}

\newtcolorbox{warnbox}{colback=warn!6,colframe=warn,boxrule=0.6pt,arc=2pt,left=6pt,right=6pt,top=4pt,bottom=4pt}
\newtcolorbox{tipbox}{colback=accent!5,colframe=accent,boxrule=0.6pt,arc=2pt,left=6pt,right=6pt,top=4pt,bottom=4pt}
\newtcolorbox{cmdbox}{colback=black!4,colframe=black!25,boxrule=0.4pt,arc=2pt,left=6pt,right=6pt,top=3pt,bottom=3pt,fontupper=\ttfamily\small}
\newcommand{\ui}[1]{\textsf{\textbf{#1}}}
\newcommand{\code}[1]{\texttt{#1}}
\newcommand{\arr}{\,$\rightarrow$\,}

\begin{document}

{\color{accent}\Huge\bfseries TrustTier: Deploying on L1 and L2}\\[4pt]
{\large Step-by-step guide: Sepolia redeploy, L1 vs L2 testnet measurements, Overleaf}\\[2pt]
{\small Repository: \url{https://github.com/gauravrai1704/TrustTier}\quad Live site: \url{https://gauravrai1704.github.io/TrustTier/}}

\vspace{6pt}
\begin{tipbox}
\textbf{What you will end up with:} TrustTier's contracts deployed on the Sepolia test network, verified on Etherscan, and the live website talking to them. Everything is \textbf{free} (Sepolia ETH has no real value). It takes about \textbf{20--30 minutes}, most of it waiting for a faucet.

\textbf{Who does what.} You can do every part yourself \emph{except Part~E}: on a personal GitHub repository, only the owner (Gaurav) can add repository secrets. Do Parts A--D, then add the secrets together with Gaurav (Part~E), then run Parts F--H. To run the workflow in Part~F you need to be a collaborator on the repository with \emph{Write} access (Gaurav: \ui{Settings}\arr\ui{Collaborators}\arr\ui{Add people}).
\end{tipbox}

\begin{tipbox}
\textbf{Already deployed before?} The contracts changed in this version (the proof now carries the model's weights hash, the tier bands adapt to volatility, and the first feature is utilization), so the old Sepolia contracts no longer match the website. If the secrets from Part~E are still in place, skip to \textbf{Part~F} and run the workflow again; nothing else is needed for the website. \textbf{Parts I--K} are new: they measure TrustTier on Sepolia (L1) and on L2 testnets for the paper, and upload the paper to Overleaf.
\end{tipbox}

\begin{warnbox}
\faExclamationTriangle\ \textbf{Never put real money anywhere in this guide.} Use a brand-new, testnet-only account. Never paste a private key into chat, email, a document or the repository.
\end{warnbox}

{\small\textbf{Checklist:}\quad \faSquare[regular]\ A account\quad \faSquare[regular]\ B $\geq$0.2 Sepolia ETH\quad \faSquare[regular]\ C RPC URL\quad \faSquare[regular]\ D Etherscan key\quad \faSquare[regular]\ E secrets added\quad \faSquare[regular]\ F run green\quad \faSquare[regular]\ G site shows Sepolia\quad \faSquare[regular]\ I L2 ETH\quad \faSquare[regular]\ J testnet runs green\quad \faSquare[regular]\ K Overleaf}

\section{Part A: Create a testnet-only account (\code{PRIVATE\_KEY})}
\begin{enumerate}
  \item Install \textbf{MetaMask} from \url{https://metamask.io} (Chrome, Firefox or Brave extension) and set it up. If you already use MetaMask, go to step 2.
  \item Create a separate account just for this: click the account name at the top\arr\ui{Add account or hardware wallet}\arr\ui{Add a new Ethereum account}\arr name it \ui{TrustTier deployer}.
  \item Show test networks: network menu (top left)\arr turn on \ui{Show test networks}\arr select \ui{Sepolia}.
  \item Copy the account's \textbf{address} (starts with \code{0x}, 42 characters). You will paste it into the faucet in Part~B.
  \item Export the \textbf{private key}: account menu (\faEllipsisV) \arr\ui{Account details}\arr\ui{Show private key}\arr enter your MetaMask password\arr copy it.
  \begin{itemize}
    \item It is 64 hexadecimal characters. If it does not start with \code{0x}, add \code{0x} in front.
    \item Keep it somewhere private. It becomes the \code{PRIVATE\_KEY} secret.
  \end{itemize}
\end{enumerate}
{\small MetaMask moves its menus between versions; if a label differs, look for the same option under the account's \ui{\faEllipsisV} menu.}

\section{Part B: Get free Sepolia test ETH}
Deployment uses about \textbf{13--15 million gas}, plus \textbf{0.01 test ETH} that funds the reward pool. At typical Sepolia gas prices that is roughly 0.03--0.15 ETH. The L1 measurement in Part~J deploys a second copy and needs about as much again. \textbf{Aim for at least 0.2 Sepolia ETH} if you will do Part~J (0.1 for the website alone). Part~I also bridges a little of it to the L2 testnets.

Use one or more of these faucets, pasting your deployer address:
\begin{itemize}
  \item \textbf{Google Cloud Web3 faucet}: search ``Google Cloud Sepolia faucet'' (on cloud.google.com: Web3, then Faucet). Sign in with Google, choose \ui{Ethereum Sepolia}. About 0.05 ETH per day.
  \item \textbf{Sepolia PoW faucet} (\url{https://sepolia-faucet.pk910.de}): no account needed. Start mining, leave the tab open for 30--60 minutes, then claim.
  \item \textbf{Alchemy, Infura or QuickNode faucets}: these usually require a small balance of real ETH on mainnet; skip them if you do not have that.
  \item Or ask a classmate or a professor to send you 0.1 Sepolia ETH; it costs them nothing real.
\end{itemize}
\textbf{Check the balance} at \url{https://sepolia.etherscan.io}: search your address and look at \ui{ETH Balance}.

\section{Part C: An RPC URL (\code{SEPOLIA\_RPC\_URL})}
This is how GitHub talks to Sepolia. Pick one:
\begin{itemize}
  \item \textbf{No signup:} \nolinkurl{https://ethereum-sepolia-rpc.publicnode.com}
  \item \textbf{More reliable (recommended):}
  \begin{enumerate}
    \item Create a free account at \url{https://www.alchemy.com}.
    \item Dashboard\arr\ui{Create new app}\arr name \ui{TrustTier}, chain \ui{Ethereum}, network \ui{Sepolia}.
    \item Open the app\arr\ui{Endpoints} / \ui{API key}\arr copy the \textbf{HTTPS} URL. It looks like \nolinkurl{https://eth-sepolia.g.alchemy.com/v2/AbC...}
  \end{enumerate}
  Infura works the same way: create an API key and copy the Sepolia HTTPS endpoint.
\end{itemize}
{\small Treat an Alchemy or Infura URL as a secret too: it contains your API key.}

\section{Part D: Etherscan API key (\code{ETHERSCAN\_API\_KEY}), optional}
This gives the contracts the green \ui{Contract Source Code Verified} tick on Etherscan, which professors can click through.
\begin{enumerate}
  \item Register at \url{https://etherscan.io} and confirm your email.
  \item Your profile\arr\ui{API Dashboard} (or \ui{API Keys})\arr\ui{Add}\arr name it \ui{TrustTier}\arr copy the key. One key also works for Sepolia.
\end{enumerate}

\section{Part E: Add the secrets to GitHub (repository owner)}
\begin{tipbox}
Only the repository owner can do this. Sit together, or share your screen, and \textbf{type the values directly into GitHub}; never send them over chat or email. Alternatively, Gaurav can create his own deployer account (Parts A--D) and add his own values.
\end{tipbox}
\begin{enumerate}
  \item Go to \url{https://github.com/gauravrai1704/TrustTier}\arr\ui{Settings} (top tab)\arr left sidebar \ui{Secrets and variables}\arr\ui{Actions}.
  \item Click \ui{New repository secret} once per row. \textbf{The names must match exactly:}
\end{enumerate}
\begin{center}
\begin{tabularx}{0.9\textwidth}{@{}lX@{}}
\toprule
\textbf{Name} & \textbf{Secret} \\
\midrule
\code{PRIVATE\_KEY} & the \code{0x...} key from Part A \\
\code{SEPOLIA\_RPC\_URL} & the URL from Part C \\
\code{ETHERSCAN\_API\_KEY} & the key from Part D (skip if you did not make one) \\
\bottomrule
\end{tabularx}
\end{center}
GitHub keeps them encrypted and never shows them again, not even to the owner.

\section{Part F: Run the deployment}
\begin{enumerate}
  \item Repository\arr\ui{Actions} tab\arr left list: \ui{Deploy to Sepolia}\arr\ui{Run workflow} (right side).
  \item Keep the defaults:
  \begin{itemize}
    \item \textbf{Branch:} \code{main}
    \item \textbf{Tier-2 verifier:} \code{real}
    \item \textbf{Use Chainlink's real ETH/USD feed:} \faCheckSquare\ (ticked)
    \item \textbf{Multisig owners:} leave \emph{empty}. If you fill this in, every later admin change needs those owners plus a timelock delay, which is awkward for a demo.
  \end{itemize}
  \item Click \ui{Run workflow} and open the run that appears. It takes about 5--10 minutes:
\end{enumerate}
\begin{center}
\begin{tabularx}{0.95\textwidth}{@{}lX@{}}
\toprule
\textbf{Step} & \textbf{What happens} \\
\midrule
Check secrets & Fails immediately if a secret name is wrong. \\
Deploy & About 15 transactions; the log prints each contract address. \\
Verify on Etherscan & Waits 30\,s, then verifies each contract. Allowed to fail without stopping the run. \\
Commit the deployment manifest & Adds \code{deployments/sepolia.json} to \code{main}. \\
Rebuild the DApp & Starts the website deploy. \\
\bottomrule
\end{tabularx}
\end{center}
When it is green, open the run's \ui{Summary}: it has a table of every contract, linked to Sepolia Etherscan.

\section{Part G: Check the live site}
\begin{enumerate}
  \item Wait about 2 minutes for \ui{Deploy DApp to GitHub Pages} to finish (Actions tab).
  \item Open \url{https://gauravrai1704.github.io/TrustTier/} and hard-refresh (\code{Ctrl+Shift+R}, or \code{Cmd+Shift+R} on a Mac).
  \item The top bar should show \ui{Sepolia \textperiodcentered{} direct}, not \ui{Simulation}. The \ui{Protocol} page lists the contracts with Etherscan links.
\end{enumerate}
\textbf{To act on it:}
\begin{itemize}
  \item \textbf{Sign in:} either \ui{Create account} in the app (then fund that new address from a faucet), or \ui{Browser wallet} with MetaMask on Sepolia (the deployer account already has ETH).
  \item \textbf{Positions page:} open a position (two transactions), then route it. Tier-0/1 claims, bond settlement and executing approvals all work from your account, and every transaction appears on Etherscan.
  \item \textbf{Tier-2 proofs and challenges} need the prover, which lives in the router (Part H). Without it the app says so.
\end{itemize}

\section{Part H (optional): Tier 2 on Sepolia, running the router on your laptop}
On your own machine (macOS or Linux; on Windows use WSL), from a clone of the repository. You need Git, Node.js~20, Python~3.11 and \code{make}. \code{make keys} builds the proving key, which can take several minutes and about 2\,GB of RAM.
\begin{cmdbox}
git clone https://github.com/gauravrai1704/TrustTier.git \&\& cd TrustTier\\
git pull\\
make install \&\& make keys\\
cd off-chain\\
export TRUSTTIER\_RPC\_URL="<your SEPOLIA\_RPC\_URL>"\\
export TRUSTTIER\_NETWORK=sepolia\\
export TRUSTTIER\_PRIVATE\_KEY="<a funded Sepolia key; the deployer key works>"\\
../.venv/bin/python -m uvicorn main:app --port 8000
\end{cmdbox}
Open the live site; it detects the router at \code{http://127.0.0.1:8000} automatically (or set it on the \ui{Protocol} page). Tier-2 decisions then use real EZKL proofs (about 10\,s each) and the real on-chain verifier.

{\small The \code{export} lines keep the key only in that terminal session. Do not write it into a file inside the repository.}


\section{Part I: Test ETH on the L2 testnets}
The paper compares Ethereum (L1) with rollups (L2). Part~J measures real transactions on each testnet, from the \textbf{same deployer address} as Part~A. Fund that address with about \textbf{0.02 ETH} on each L2 testnet you want (one is enough; all three is best):
\begin{center}
\renewcommand{\arraystretch}{1.2}
\begin{tabularx}{\textwidth}{@{}lX@{}}
\toprule
\textbf{Network} & \textbf{How to get test ETH there} \\
\midrule
\textbf{Base Sepolia} (recommended first) & Bridge Sepolia ETH with the official testnet bridge \url{https://testnets.superbridge.app} (connect MetaMask, choose \ui{Sepolia}\arr\ui{Base Sepolia}, about 3--10 minutes), or use a Base Sepolia faucet (Coinbase Developer Platform, Alchemy). \\
\textbf{OP Sepolia} & Superchain faucet \url{https://console.optimism.io/faucet}, or bridge from Sepolia at \url{https://testnets.superbridge.app}. \\
\textbf{Arbitrum Sepolia} & Bridge from Sepolia at \url{https://bridge.arbitrum.io} (switch on testnet mode), or an Arbitrum Sepolia faucet (Alchemy, QuickNode). \\
\bottomrule
\end{tabularx}
\end{center}
In MetaMask, add the network if it is missing: network menu\arr\ui{Add a custom network}, or open \url{https://chainlist.org}, tick \ui{Include testnets}, search the name and click \ui{Add to MetaMask}. Check the balance on \url{https://sepolia.basescan.org}, \url{https://sepolia-optimism.etherscan.io} or \url{https://sepolia.arbiscan.io}.

{\small Nothing else is needed: the workflow uses public RPCs for the L2 testnets. If one is slow, you can add your own endpoint as a secret (Part~E): \code{BASE\_SEPOLIA\_RPC\_URL}, \code{OP\_SEPOLIA\_RPC\_URL} or \code{ARBITRUM\_SEPOLIA\_RPC\_URL} (Alchemy offers all three).}

\section{Part J: Measure real decisions on L1 and L2 (for the paper)}
This deploys a separate \emph{benchmark} copy of TrustTier (it does not touch the website's contracts), sends one of each decision with the real verifier and a real zero-knowledge proof, runs the on-chain-inference baseline next to it, and commits what every transaction actually paid, including the fee a rollup charges for posting data to Ethereum.
\begin{enumerate}
  \item Repository\arr\ui{Actions}\arr left list: \ui{Testnet costs (L1 vs L2)}\arr\ui{Run workflow}.
  \item \textbf{Branch:} \code{main}. \textbf{Network:} pick one of \code{sepolia}, \code{baseSepolia}, \code{opSepolia}, \code{arbitrumSepolia}.
  \item Click \ui{Run workflow}. It takes about 5--10 minutes. Then run it again for each other network you funded. Run \code{sepolia} at least once: it is the L1 side of the comparison.
\end{enumerate}
\begin{center}
\begin{tabularx}{0.95\textwidth}{@{}lX@{}}
\toprule
\textbf{Step} & \textbf{What happens} \\
\midrule
Deploy a benchmark copy & About 20 transactions; writes \path{deployments/<network>-bench.json}. \\
Send real decisions & Tier~0, Tier~1 + reveal, Tier~2 with a real proof, the baseline's deploy, weights and decision, and a liquidation. The log prints each transaction's gas and fee. \\
Commit the results & Adds \path{off-chain/evaluation/results/testnet_costs/<network>.json} to \code{main}. \\
\bottomrule
\end{tabularx}
\end{center}
The run's \ui{Summary} shows a table of every action with its gas and fee. When the runs you want are green, tell the team (or Claude): the numbers go into the paper's L1/L2 section, next to the live-mainnet fee study it already has.

\section{Part K: Upload the paper to Overleaf}
\begin{enumerate}
  \item Download \code{TrustTier\_overleaf.zip} (it holds \code{main.tex}, \code{references.bib} and the \code{figures/} folder with every figure pre-rendered).
  \item On \url{https://www.overleaf.com}: \ui{New Project}\arr\ui{Upload Project}\arr choose the zip.
  \item Open \code{main.tex} and click \ui{Recompile}. The default settings work (pdfLaTeX, BibTeX). Figures are pre-rendered PDFs in \path{figures/ext/}, so it compiles within the free plan's time limit.
  \item To change a figure: edit its TikZ code, delete that figure's PDF in \path{figures/ext/}, and enable shell escape (\ui{Menu}\arr\ui{Settings}; Overleaf allows it) for one compile to re-render it.
\end{enumerate}

\section{If something goes wrong}
\begin{center}
\renewcommand{\arraystretch}{1.25}
\begin{tabularx}{\textwidth}{@{}>{\raggedright\arraybackslash}p{6.2cm}X@{}}
\toprule
\textbf{What you see} & \textbf{Fix} \\
\midrule
``Add the SEPOLIA\_RPC\_URL and PRIVATE\_KEY repository secrets first'' & A secret is missing or its name is misspelled; recheck Part E. \\
\code{insufficient funds} & Get more test ETH (Part B), then click \ui{Re-run all jobs}. \\
\code{429} / rate limit / timeout from the RPC & The public node is busy: use an Alchemy URL (Part C) and re-run. \\
\code{replacement transaction underpriced} / nonce errors & Wait 2 minutes for pending transactions, then re-run. \\
Verify step shows ``could not verify'' & Harmless; usually a missing or brand-new Etherscan key. Verify later from a laptop with \code{ETHERSCAN\_API\_KEY} set: \texttt{npx hardhat run scripts/verify.js -{}-network sepolia}. \\
Commit step fails (``protected branch'') & Owner: \ui{Settings}\arr\ui{Branches}: allow GitHub Actions to push to \code{main}, or remove the protection. \\
Site still says ``Simulation'' & Wait for the Pages deploy to finish, then hard-refresh. \\
\ui{Run workflow} button missing & You need Write access to the repository (see the box on page 1). \\
Testnet costs: \code{insufficient funds} on an L2 & The deployer has no ETH on that L2 yet; bridge or use a faucet (Part~I), then re-run. \\
Testnet costs: \code{features ... differ from the proof's} & A feed or the clock moved between steps (rare on a busy public RPC). Re-run; if it repeats, add your own RPC secret for that network. \\
Testnet costs: commit step fails & Another run pushed first; click \ui{Re-run failed jobs}. \\
\bottomrule
\end{tabularx}
\end{center}

\section{Good to know}
\begin{itemize}
  \item \textbf{Re-running the deploy} creates a fresh set of contracts and overwrites \path{deployments/sepolia.json}; the site then switches to the new ones.
  \item The \ui{Refresh Sepolia demo feeds} workflow runs automatically every 12 hours with the same secrets. It spends a tiny amount of test ETH and keeps the demo's mock price feeds from going stale before the presentation. Turn it off under \ui{Actions}\arr that workflow\arr\ui{\faEllipsisH}\arr\ui{Disable workflow}.
  \item \textbf{Never reuse this key for anything with real money.} If you think it leaked, make a new account and have the owner update the \code{PRIVATE\_KEY} secret.
  \item The benchmark copies from Part~J are separate from the website's contracts; they can be ignored afterwards.
  \item When you are done, send the link to the workflow run (or the error from its log) to the team so the deployment and the live site can be checked.
\end{itemize}

\end{document}
